% =============================================================================
%  HAR-RV: methodological specification (raw and logarithmic)
%
%  Compiles standalone:      pdflatex har_rv_methodology.tex
%  Or drop into a thesis:    copy everything between the two
%                            "%%% BODY BEGIN/END" markers into a chapter and
%                            move the \newcommand block to your preamble.
%
%  Companion prose version:  docs/HAR-RV-methodology.md
%  Reference implementation: HAR-RV_RUN.PY
% =============================================================================
\documentclass[11pt,a4paper]{article}

\usepackage[T1]{fontenc}
\usepackage[utf8]{inputenc}
\usepackage{amsmath,amssymb,amsthm}
\usepackage{booktabs}
\usepackage{bm}
\usepackage[margin=2.6cm]{geometry}
\usepackage[hidelinks]{hyperref}

% ---- notation shorthands (move to preamble if \input-ing this file) ---------
\newcommand{\RV}{\mathrm{RV}}
\newcommand{\BPV}{\mathrm{BPV}}
\newcommand{\RQ}{\mathrm{RQ}}
\newcommand{\IV}{\mathrm{IV}}
\newcommand{\QV}{\mathrm{QV}}
\newcommand{\E}{\mathbb{E}}
\newcommand{\Var}{\operatorname{Var}}
\newcommand{\Cov}{\operatorname{Cov}}
\newcommand{\Fset}{\mathcal{F}}
\newcommand{\MSE}{\mathrm{MSE}}
\newcommand{\MAE}{\mathrm{MAE}}
\newcommand{\QLIKE}{\mathrm{QLIKE}}
\newcommand{\dd}{\mathrm{d}}

\title{The HAR-RV Model:\\ Methodological Specification in Levels and Logarithms}
\date{}

\begin{document}
\maketitle

%%% BODY BEGIN %%%

\section{Notation and timing}
\label{sec:notation}

Let $t=1,\dots,T$ index trading days and let the log price within day $t$ be sampled on an
intraday grid of $M$ intervals, $r_{t,j}=p_{t,j}-p_{t,j-1}$ for $j=1,\dots,M$. Write
$\Fset_t$ for the information set generated by all data observed up to and including the
close of day~$t$.

Throughout, regressors are dated $t$ and measured on $\Fset_t$, and the dependent variable
spans days $t+1,\dots,t+h$; no contemporaneous information enters the right-hand side. (The
accompanying code indexes each row by the \emph{first day of the forecast window} rather
than by the last day of the information set; substituting $t\to t+1$ maps one convention
into the other and leaves every regression unchanged.)

\section{The volatility measure}
\label{sec:measure}

Daily realized variance is
\begin{equation}
  \RV_t \;=\; \sum_{j=1}^{M} r_{t,j}^{2}.
  \label{eq:rv}
\end{equation}
For a semimartingale log price $\dd p_s=\mu_s\,\dd s+\sigma_s\,\dd W_s+\dd J_s$,
\begin{equation}
  \RV_t \;\xrightarrow[M\to\infty]{p}\;
  \underbrace{\int_{t-1}^{t}\sigma_s^{2}\,\dd s}_{\IV_t}
  \;+\;\sum_{t-1<s\le t}\bigl(\Delta J_s\bigr)^{2}
  \;=\;\QV_t ,
  \label{eq:qv}
\end{equation}
so $\RV_t$ consistently estimates quadratic variation, which coincides with integrated
variance in the absence of jumps. The model is specified on the \emph{variance}, not the
standard deviation; units must be stated, since loss levels (though not loss ratios, QLIKE
or $R^{2}$) depend on them. Observations with $\RV_t\le 0$ are non-trading days and are
removed: $\ln \RV_t$ and QLIKE are otherwise undefined.

\section{Heterogeneous components}
\label{sec:components}

The daily, weekly ($5$-day) and monthly ($22$-day) components are backward-looking
arithmetic averages,
\begin{align}
  \RV^{(d)}_t &= \RV_t, \label{eq:comp-d}\\[2pt]
  \RV^{(w)}_t &= \frac{1}{5}\sum_{i=0}^{4}\RV_{t-i}, \label{eq:comp-w}\\[2pt]
  \RV^{(m)}_t &= \frac{1}{22}\sum_{i=0}^{21}\RV_{t-i}. \label{eq:comp-m}
\end{align}
The windows are nested, so the three regressors are strongly collinear and the coefficients
are \emph{partial} effects conditional on the remaining averages; a correlation matrix and
variance-inflation factors should accompany any coefficient table. The lengths $5$ and $22$
are trading-day conventions, not estimated quantities.

Equations \eqref{eq:comp-d}--\eqref{eq:comp-m} approximate the partial volatility cascade of
Corsi (2009), in which agents operating at monthly, weekly and daily horizons form
expectations that feed downward:
\begin{align}
  \sigma^{(m)}_{t+1m} &= c^{(m)}+\phi^{(m)}\RV^{(m)}_{t}+\tilde\omega^{(m)}_{t+1m},\notag\\
  \sigma^{(w)}_{t+1w} &= c^{(w)}+\phi^{(w)}\RV^{(w)}_{t}
                         +\gamma^{(w)}\E_t\bigl[\sigma^{(m)}_{t+1m}\bigr]+\tilde\omega^{(w)}_{t+1w},
                         \label{eq:cascade}\\
  \sigma^{(d)}_{t+1d} &= c^{(d)}+\phi^{(d)}\RV^{(d)}_{t}
                         +\gamma^{(d)}\E_t\bigl[\sigma^{(w)}_{t+1w}\bigr]+\tilde\omega^{(d)}_{t+1d}.\notag
\end{align}
Recursive substitution, together with $\RV_{t+1}=\sigma^{(d)}_{t+1d}+\omega_{t+1}$, delivers
the three-factor structure that the regressions below approximate. The model is not a formal
long-memory process: it superposes three exponentially decaying components to mimic
hyperbolic autocorrelation decay.

\section{Dependent variable}
\label{sec:target}

For horizon $h\in\{1,5,22\}$ the $h$-day forward \emph{average} realized variance is
\begin{equation}
  \RV_{t,\,t+h} \;=\; \frac{1}{h}\sum_{i=1}^{h}\RV_{t+i},
  \qquad \RV_{t,\,t+1}=\RV_{t+1}.
  \label{eq:target}
\end{equation}
Using the $h$-day sum instead scales the dependent variable, all coefficients and all
residuals by $h$, and the MSE by $h^{2}$; $R^{2}$ is invariant but reported losses are not,
and horizons cease to be comparable.

\section{Specification I: levels}
\label{sec:raw}

\begin{equation}
  \boxed{\;
  \RV_{t,\,t+h}
  \;=\;\beta_{0}
  +\beta_{d}\,\RV^{(d)}_{t}
  +\beta_{w}\,\RV^{(w)}_{t}
  +\beta_{m}\,\RV^{(m)}_{t}
  +\varepsilon^{(h)}_{t+h},
  \qquad t=22,\dots,T-h,\;}
  \label{eq:har-raw}
\end{equation}
with $\E\bigl[\varepsilon^{(h)}_{t+h}\mid\Fset_t\bigr]=0$. The monthly component requires
$\RV_{t-21}$ and the target requires $\RV_{t+h}$, so the effective sample size is
\begin{equation}
  T_h \;=\; T-h-21,
  \label{eq:Th}
\end{equation}
which differs across horizons and should be reported per horizon rather than as the raw
series length $T$.

In levels the residuals are strongly heteroskedastic (dispersion scales with the level of
$\RV$), right-skewed and leptokurtic, the fit is dominated by a small number of crisis
observations, and fitted values are not constrained to be positive.

\section{Specification II: logarithms}
\label{sec:log}

\subsection{Log-of-mean convention}

The logarithm is applied \emph{outside} every sum:
\begin{equation}
  x^{(d)}_t=\ln \RV_t,\qquad
  x^{(w)}_t=\ln\!\Bigl(\tfrac{1}{5}\textstyle\sum_{i=0}^{4}\RV_{t-i}\Bigr),\qquad
  x^{(m)}_t=\ln\!\Bigl(\tfrac{1}{22}\textstyle\sum_{i=0}^{21}\RV_{t-i}\Bigr),
  \label{eq:log-regressors}
\end{equation}
\begin{equation}
  y_{t,\,t+h}\;=\;\ln \RV_{t,\,t+h}\;=\;\ln\!\Bigl(\tfrac{1}{h}\textstyle\sum_{i=1}^{h}\RV_{t+i}\Bigr),
  \label{eq:log-target}
\end{equation}
\begin{equation}
  \boxed{\;
  y_{t,\,t+h}
  \;=\;\beta_{0}
  +\beta_{d}\,x^{(d)}_{t}
  +\beta_{w}\,x^{(w)}_{t}
  +\beta_{m}\,x^{(m)}_{t}
  +u^{(h)}_{t+h},
  \qquad t=22,\dots,T-h.\;}
  \label{eq:har-log}
\end{equation}
Equivalently $x^{(\cdot)}_t=\ln \RV^{(\cdot)}_t$ and $y_{t,t+h}=\ln \RV_{t,t+h}$: coefficients
are elasticities, a one-percent rise in a component raising the expected $h$-day forward
average by $\beta_{\cdot}$ percent.

\subsection{Why the placement of the logarithm matters}

By Jensen's inequality, for any window of length $n$,
\begin{equation}
  \underbrace{\frac{1}{n}\sum_{i=0}^{n-1}\ln \RV_{t-i}}_{\ln(\text{geometric mean})}
  \;\le\;
  \underbrace{\ln\!\Bigl(\frac{1}{n}\sum_{i=0}^{n-1}\RV_{t-i}\Bigr)}_{\ln(\text{arithmetic mean})},
  \label{eq:jensen}
\end{equation}
with equality only if $\RV$ is constant across the window; the gap is widest exactly when
$\RV$ is most dispersed. Both conventions appear in the literature. Log-of-mean is adopted
here because it delivers the exact identity
\begin{equation}
  \exp\bigl(y_{t,\,t+h}\bigr)\;=\;\RV_{t,\,t+h}\qquad\text{for every }h,
  \label{eq:identity}
\end{equation}
so that \eqref{eq:har-raw} and \eqref{eq:har-log} target the same object at every horizon and
back-transformed losses are directly comparable. Under mean-of-logs, $\exp(y)$ would be a
\emph{geometric} forward mean---a smaller, mechanically smoother quantity---and cross-scale
comparisons would be legitimate only at $h=1$. At $h=1$, and for the daily component, the
two conventions coincide, which makes $h=1$ a convenient unit test.

Sums versus means are cosmetic on the log scale, since
$\ln\bigl(\sum_i \RV_{t-i}\bigr)=\ln\bigl(n^{-1}\sum_i \RV_{t-i}\bigr)+\ln n$ and the constant
is absorbed by the intercept.

\subsection{Back-transformation}

Since $\exp(\cdot)$ is convex, $\exp\bigl(\E[y\mid\Fset_t]\bigr)$ is the conditional median
and understates the conditional mean. Under conditional normality of $u^{(h)}_{t+h}$,
\begin{equation}
  \E\bigl[\RV_{t,\,t+h}\mid\Fset_t\bigr]
  =\exp\!\Bigl(\E\bigl[y_{t,\,t+h}\mid\Fset_t\bigr]+\tfrac{1}{2}\sigma_u^{2}\Bigr),
  \label{eq:jensen-correction}
\end{equation}
implemented as
\begin{equation}
  \widehat{\RV}_{t,\,t+h}
  =\exp\!\Bigl(\hat y_{t,\,t+h}+\tfrac{1}{2}\hat\sigma_u^{2}\Bigr),
  \qquad
  \hat\sigma_u^{2}=\frac{1}{T_h-K}\sum_{t}\hat u_t^{2},
  \label{eq:backtransform}
\end{equation}
with $\hat\sigma_u^{2}$ the in-sample residual variance and $K=4$. The correction is a
systematic level shift, not noise, and does not average out. If normality is doubtful,
Duan's (1983) smearing estimator
$\widehat{\RV}=\exp(\hat y)\,T_h^{-1}\sum_t\exp(\hat u_t)$ replaces $\exp(\hat\sigma_u^{2}/2)$.
By construction $\widehat{\RV}_{t,t+h}>0$, so the log specification cannot produce the
negative variance forecasts that \eqref{eq:har-raw} can.

\section{Estimation and inference}
\label{sec:estimation}

\subsection{OLS}

With $z_t=\bigl(1,\RV^{(d)}_t,\RV^{(w)}_t,\RV^{(m)}_t\bigr)'$ in levels, or
$z_t=\bigl(1,x^{(d)}_t,x^{(w)}_t,x^{(m)}_t\bigr)'$ in logs, $K=4$, and design matrix $Z$,
\begin{equation}
  \hat\beta^{(h)}=\bigl(Z'Z\bigr)^{-1}Z'y
  =\Bigl(\sum_t z_t z_t'\Bigr)^{-1}\Bigl(\sum_t z_t\,y_{t,\,t+h}\Bigr).
  \label{eq:ols}
\end{equation}
A separate regression is estimated for each $h$ (direct multi-step projection), rather than
iterating the $h=1$ model forward.

\subsection{Overlapping observations}

For $h>1$ successive dependent variables share $h-1$ daily terms. If the one-day innovations
$\nu_t$ are serially uncorrelated, then
$\varepsilon^{(h)}_{t+h}\approx h^{-1}\sum_{i=1}^{h}\nu_{t+i}$ and
\begin{equation}
  \Cov\bigl(\varepsilon^{(h)}_{t+h},\,\varepsilon^{(h)}_{t+h+j}\bigr)
  =\frac{h-|j|}{h^{2}}\,\sigma_\nu^{2}\,\mathbf{1}\{|j|<h\},
  \label{eq:ma}
\end{equation}
an $\mathrm{MA}(h-1)$ error. OLS remains consistent, since
$\E[z_t\varepsilon^{(h)}_{t+h}]=0$ is unaffected, but conventional standard errors are
severely understated: autocorrelation-robust inference is mandatory at $h=5$ and $h=22$.

\subsection{Newey--West HAC covariance}

With residuals $\hat e_t$, Bartlett weights $w_j=1-j/(L+1)$ and bandwidth $L$,
\begin{equation}
  \hat S_L=\sum_{t}\hat e_t^{2}z_tz_t'
  +\sum_{j=1}^{L}w_j\sum_{t=j+1}\hat e_t\hat e_{t-j}\bigl(z_tz_{t-j}'+z_{t-j}z_t'\bigr),
  \label{eq:hac-S}
\end{equation}
\begin{equation}
  \widehat{\Var}\bigl(\hat\beta^{(h)}\bigr)
  =\frac{T_h}{T_h-K}\bigl(Z'Z\bigr)^{-1}\hat S_L\bigl(Z'Z\bigr)^{-1},
  \label{eq:hac-V}
\end{equation}
the leading factor being the finite-sample degrees-of-freedom correction. Bartlett weights
ensure $\hat S_L$ is positive semi-definite. All reported $t$-statistics use
\eqref{eq:hac-V}.

\subsection{Bandwidth}

\begin{equation}
  L(h,T_h)=\max\Bigl\{\,\underbrace{2(h-1)}_{\text{overlap}},\;
  \underbrace{\bigl\lfloor 4\,(T_h/100)^{2/9}\bigr\rfloor}_{\text{Newey--West (1994) plug-in}}\Bigr\}.
  \label{eq:bandwidth}
\end{equation}
The first term covers autocorrelation induced mechanically by the overlapping target;
\eqref{eq:ma} implies an $\mathrm{MA}(h-1)$ error, so $h-1$ is the theoretical minimum and
the factor of two is the conservative convention standard in this literature. The second is
data-driven and addresses serial correlation present in the series itself. The maximum
matters at $h=1$, where $2(h-1)=0$ would assert serially uncorrelated residuals; a Bartlett
kernel with $L=0$ retains only the $j=0$ term of \eqref{eq:hac-S} and reduces to White/HC1,
which is heteroskedasticity-robust but not a HAC correction. Report $L$ for each horizon,
and the pure $2(h-1)$ rule as a robustness check.

\section{Evaluation}
\label{sec:evaluation}

\subsection{Sample split}

The split is chronological, never random:
\begin{equation}
  \mathcal{T}_{\text{train}}=\{t:\text{year}(t)\le 2024\},\qquad
  \mathcal{T}_{\text{test}}=\{t:\text{year}(t)\ge 2025\}.
  \label{eq:split}
\end{equation}
Coefficients are estimated once on $\mathcal{T}_{\text{train}}$ and held fixed over
$\mathcal{T}_{\text{test}}$ (fixed scheme). OLS has no hyperparameters, so no validation
split is needed; when benchmarking against models that require one, their validation years
may be folded into the HAR estimation sample provided the test window is identical. Two
caveats deserve a sentence: parameters become progressively stale across a long test window
(a rolling-window re-estimation is the standard robustness check), and $h-1$ observations
straddle the train/test boundary.

\subsection{Loss functions}

Over $n$ test observations, with $\hat\sigma^{2}_t$ the forecast and $\RV_t$ the proxy for
the latent conditional variance,
\begin{equation}
  \MSE=\frac{1}{n}\sum_t\bigl(\RV_t-\hat\sigma^{2}_t\bigr)^{2},
  \qquad
  \MAE=\frac{1}{n}\sum_t\bigl|\RV_t-\hat\sigma^{2}_t\bigr|,
  \label{eq:mse-mae}
\end{equation}
\begin{equation}
  \QLIKE=\frac{1}{n}\sum_t\left[\frac{\RV_t}{\hat\sigma^{2}_t}
  -\ln\!\left(\frac{\RV_t}{\hat\sigma^{2}_t}\right)-1\right]\;\ge\;0,
  \label{eq:qlike}
\end{equation}
QLIKE being written in normalised form, so that it vanishes only for a perfect forecast and
penalises under-prediction of volatility more heavily than over-prediction.

Patton (2011) shows that MSE and QLIKE are \emph{robust} to noise in the volatility proxy:
because $\RV_t$ is conditionally unbiased for $\sigma^{2}_t$, forecast rankings under the
proxy coincide with rankings under the latent target. MAE does not have this property and
can reverse the true ranking; it may be reported descriptively, but model selection should
rest on MSE and QLIKE.

In levels, all three losses are computed on the $\RV$ scale. In logs, MSE and MAE are
reported in $\ln \RV$ units---the scale on which the model is fitted, where squared error is
a relative-error criterion---while QLIKE requires variances and is evaluated after
back-transformation by \eqref{eq:backtransform}. Identity \eqref{eq:identity} makes
back-transformed losses comparable with the levels specification at every horizon. Reporting
QLIKE both with and without the Jensen term makes the size of the correction visible.

Levels OLS can return $\hat\sigma^{2}_t\le 0$, where QLIKE is undefined. Such forecasts are
floored at $\varsigma\,\overline{\RV}_{\text{train}}$ with $\varsigma=10^{-4}$ and the count
is reported; discarding them would delete the worst forecasts and bias QLIKE downward. The
floor never binds under \eqref{eq:backtransform}.

\subsection{Forecast comparison}

For two competing forecasts with loss differential
$d_t=L(\RV_t,\hat\sigma^{2}_{1,t})-L(\RV_t,\hat\sigma^{2}_{2,t})$, the Diebold--Mariano
(1995) statistic is
\begin{equation}
  \mathrm{DM}=\frac{\bar d}{\sqrt{\widehat{\Var}_{\mathrm{HAC}}(\bar d)}}
  \;\xrightarrow{d}\;\mathcal{N}(0,1),
  \qquad \bar d=\frac{1}{n}\sum_t d_t,
  \label{eq:dm}
\end{equation}
with the HAC variance computed under bandwidth rule \eqref{eq:bandwidth}, since the loss
differential inherits the same overlap. With more than two models, the Model Confidence Set
of Hansen, Lunde and Nason (2011) controls the multiple-comparison problem.

\section{Residual diagnostics}
\label{sec:diagnostics}

Computed on training residuals for each horizon and specification.

\begin{center}
\begin{tabular}{@{}llp{4.4cm}@{}}
\toprule
Test & Statistic & Null hypothesis \\
\midrule
Durbin--Watson  & $DW=\sum_{t\ge2}(\hat e_t-\hat e_{t-1})^{2}\big/\sum_t\hat e_t^{2}$ & no first-order autocorrelation \\
Ljung--Box      & $Q(L)=T_h(T_h+2)\sum_{j=1}^{L}\hat\rho_j^{2}/(T_h-j)\sim\chi^{2}_{L}$ & no autocorrelation to lag $L$ \\
Jarque--Bera    & $JB=\tfrac{T_h}{6}\bigl(S^{2}+\tfrac{1}{4}(\kappa-3)^{2}\bigr)\sim\chi^{2}_{2}$ & normality \\
Breusch--Pagan  & $LM=T_h R^{2}_{\text{aux}}\sim\chi^{2}_{K-1}$ & homoskedasticity \\
\bottomrule
\end{tabular}
\end{center}

These are not a pass/fail gate but the justification for the HAC covariance and for the log
transform. A large Breusch--Pagan statistic in levels is expected; it confirms that
classical standard errors would be understated. Logging reduces heteroskedasticity and
non-normality sharply---it converts multiplicative volatility-of-volatility into
approximately additive noise---but does not remove serial correlation, so
\eqref{eq:hac-V} remains necessary in both specifications.

\section{Forecast construction}
\label{sec:forecast}

\begin{align}
  \widehat{\RV}_{t,\,t+h}
  &=\hat\beta_{0}+\hat\beta_{d}\RV^{(d)}_{t}+\hat\beta_{w}\RV^{(w)}_{t}+\hat\beta_{m}\RV^{(m)}_{t}
  && \text{(levels)}, \label{eq:fc-raw}\\[2pt]
  \hat y_{t,\,t+h}
  &=\hat\beta_{0}+\hat\beta_{d}x^{(d)}_{t}+\hat\beta_{w}x^{(w)}_{t}+\hat\beta_{m}x^{(m)}_{t},
  \quad
  \widehat{\RV}_{t,\,t+h}=\exp\!\Bigl(\hat y_{t,\,t+h}+\tfrac{1}{2}\hat\sigma_u^{2}\Bigr)
  && \text{(logs)}. \label{eq:fc-log}
\end{align}
If annualised volatility is the reported quantity, apply
$\sqrt{252\,\widehat{\RV}_{t,t+h}}$ \emph{after} the Jensen correction, noting that the
square root introduces its own, opposite, concavity gap.

\section{Analytical properties}
\label{sec:properties}

\subsection{Implied constrained AR(22)}

At $h=1$ the levels HAR is an AR(22) with only three free slope parameters. Writing
$\RV_{t+1}=\beta_0+\sum_{j=1}^{22}\phi_j\RV_{t-j+1}+\varepsilon_{t+1}$,
\begin{equation}
  \phi_1=\beta_d+\frac{\beta_w}{5}+\frac{\beta_m}{22},
  \qquad
  \phi_j=\frac{\beta_w}{5}+\frac{\beta_m}{22}\;\;(j=2,\dots,5),
  \qquad
  \phi_j=\frac{\beta_m}{22}\;\;(j=6,\dots,22).
  \label{eq:ar22}
\end{equation}
The step-function shape of \eqref{eq:ar22} is what generates slowly decaying,
quasi-hyperbolic autocorrelation from three parameters.

\subsection{Persistence and stationarity}

\begin{equation}
  \sum_{j=1}^{22}\phi_j=\beta_d+\beta_w+\beta_m,
  \qquad
  \E[\RV]=\frac{\beta_0}{1-\beta_d-\beta_w-\beta_m}.
  \label{eq:persistence}
\end{equation}
Covariance stationarity requires the roots of $1-\sum_{j=1}^{22}\phi_jL^{j}$ to lie outside
the unit circle; when $\beta_d,\beta_w,\beta_m\ge0$---so that all $\phi_j\ge0$---this reduces
to $\beta_d+\beta_w+\beta_m<1$. Estimates in the range $0.95$--$0.99$ are typical. Report the
sum with its HAC standard error, from $c'\widehat{\Var}(\hat\beta)c$ with $c=(0,1,1,1)'$.

\subsection{Scale equivariance}

Under $\RV\to c\,\RV$ with $c>0$: in levels $\beta_0\to c\beta_0$ with slopes unchanged, MSE
scaling by $c^{2}$ and MAE by $c$, while QLIKE and $R^{2}$ are invariant; in logs
$\beta_0\to\beta_0+(1-\beta_d-\beta_w-\beta_m)\ln c$ with slopes unchanged and log-scale
losses invariant. Loss \emph{levels} are therefore incomparable across studies unless units
are stated; loss ratios and QLIKE are comparable.

\section{Standard extensions}
\label{sec:extensions}

Each retains the estimation and inference machinery of Section~\ref{sec:estimation} and
changes only the regressor set. With bipower variation
$\BPV_t=\tfrac{\pi}{2}\tfrac{M}{M-1}\sum_{j=2}^{M}|r_{t,j}||r_{t,j-1}|$, a jump-robust
estimator of $\IV_t$, the jump measure is $J_t=\max(\RV_t-\BPV_t,0)$ and the HAR-J is
\begin{equation}
  \RV_{t,t+h}=\beta_0+\beta_d\RV^{(d)}_t+\beta_w\RV^{(w)}_t+\beta_m\RV^{(m)}_t
  +\beta_j J_t+\varepsilon^{(h)}_{t+h}.
  \label{eq:harj}
\end{equation}
HAR-CJ replaces $\RV$ throughout by its continuous part, with jumps identified by a
significance test rather than truncation at zero. LHAR adds negative-return aggregates to
capture leverage. SHAR splits $\RV$ into signed semivariances
$\RV^{\pm}_t=\sum_j r_{t,j}^{2}\mathbf{1}\{\pm r_{t,j}>0\}$, with $\RV^{+}_t+\RV^{-}_t=\RV_t$.
HARQ makes the daily coefficient depend on the measurement error of $\RV_t$ through realized
quarticity $\RQ_t=\tfrac{M}{3}\sum_{j=1}^{M}r_{t,j}^{4}$:
\begin{equation}
  \RV_{t,t+h}=\beta_0+\bigl(\beta_d+\beta_{d,Q}\RQ_t^{1/2}\bigr)\RV^{(d)}_t
  +\beta_w\RV^{(w)}_t+\beta_m\RV^{(m)}_t+\varepsilon^{(h)}_{t+h}.
  \label{eq:harq}
\end{equation}

\section{Assumptions}
\label{sec:assumptions}

\begin{enumerate}
  \item \textbf{Proxy consistency.} $\RV_t$ is a consistent, conditionally unbiased estimator
        of $\QV_t$; measurement error from finite $M$, microstructure noise and
        non-synchronous trading is ignored in \eqref{eq:har-raw} and \eqref{eq:har-log}.
        HARQ \eqref{eq:harq} is the standard response.
  \item \textbf{Linearity.} The conditional mean is linear in the three components---an
        approximation to the cascade \eqref{eq:cascade}, not an implication of it.
  \item \textbf{Stationarity and weak dependence.} $\{\RV_t\}$ is covariance-stationary and
        mixing, delivering consistency and asymptotic normality under \eqref{eq:hac-V}.
        Realized volatility is highly persistent, so this is the assumption most exposed to
        challenge; report ADF or KPSS tests on $\RV_t$ and $\ln \RV_t$, noting their low
        power against near-long-memory alternatives.
  \item \textbf{Parameter constancy} over the estimation window, per the fixed scheme
        \eqref{eq:split}; structural breaks around crises are common.
  \item \textbf{Conditional log-normality}, required only for \eqref{eq:jensen-correction},
        not for estimation or for log-scale losses.
\end{enumerate}

%%% BODY END %%%

\begin{thebibliography}{99}
\bibitem{ab1998} Andersen, T.~G., \& Bollerslev, T. (1998). Answering the skeptics: Yes, standard volatility models do provide accurate forecasts. \emph{International Economic Review}, 39(4), 885--905.
\bibitem{abd2007} Andersen, T.~G., Bollerslev, T., \& Diebold, F.~X. (2007). Roughing it up: Including jump components in the measurement, modeling, and forecasting of return volatility. \emph{Review of Economics and Statistics}, 89(4), 701--720.
\bibitem{abdl2003} Andersen, T.~G., Bollerslev, T., Diebold, F.~X., \& Labys, P. (2003). Modeling and forecasting realized volatility. \emph{Econometrica}, 71(2), 579--625.
\bibitem{bns2004} Barndorff-Nielsen, O.~E., \& Shephard, N. (2004). Power and bipower variation with stochastic volatility and jumps. \emph{Journal of Financial Econometrics}, 2(1), 1--37.
\bibitem{bpq2016} Bollerslev, T., Patton, A.~J., \& Quaedvlieg, R. (2016). Exploiting the errors: A simple approach for improved volatility forecasting. \emph{Journal of Econometrics}, 192(1), 1--18.
\bibitem{corsi2009} Corsi, F. (2009). A simple approximate long-memory model of realized volatility. \emph{Journal of Financial Econometrics}, 7(2), 174--196.
\bibitem{cr2012} Corsi, F., \& Ren\`o, R. (2012). Discrete-time volatility forecasting with persistent leverage effect and the link with continuous-time volatility modeling. \emph{Journal of Business \& Economic Statistics}, 30(3), 368--380.
\bibitem{dm1995} Diebold, F.~X., \& Mariano, R.~S. (1995). Comparing predictive accuracy. \emph{Journal of Business \& Economic Statistics}, 13(3), 253--263.
\bibitem{duan1983} Duan, N. (1983). Smearing estimate: A nonparametric retransformation method. \emph{Journal of the American Statistical Association}, 78(383), 605--610.
\bibitem{hln2011} Hansen, P.~R., Lunde, A., \& Nason, J.~M. (2011). The model confidence set. \emph{Econometrica}, 79(2), 453--497.
\bibitem{nw1987} Newey, W.~K., \& West, K.~D. (1987). A simple, positive semi-definite, heteroskedasticity and autocorrelation consistent covariance matrix. \emph{Econometrica}, 55(3), 703--708.
\bibitem{nw1994} Newey, W.~K., \& West, K.~D. (1994). Automatic lag selection in covariance matrix estimation. \emph{Review of Economic Studies}, 61(4), 631--653.
\bibitem{patton2011} Patton, A.~J. (2011). Volatility forecast comparison using imperfect volatility proxies. \emph{Journal of Econometrics}, 160(1), 246--256.
\bibitem{ps2009} Patton, A.~J., \& Sheppard, K. (2009). Evaluating volatility and correlation forecasts. In T.~G. Andersen et al. (Eds.), \emph{Handbook of Financial Time Series} (pp.~801--838). Springer.
\bibitem{ps2015} Patton, A.~J., \& Sheppard, K. (2015). Good volatility, bad volatility: Signed jumps and the persistence of volatility. \emph{Review of Economics and Statistics}, 97(3), 683--697.
\end{thebibliography}

\end{document}
